\documentclass[12pt,a4paper]{article}

% --- PAQUETES FUNDAMENTALES PARA OVERLEAF ---
\usepackage[utf8]{inputenc}
\usepackage[spanish,es-nodecimaldot]{babel}
\usepackage[margin=2.5cm]{geometry}
\usepackage{graphicx}
\usepackage{amsmath,amssymb}
\usepackage{booktabs}
\usepackage{array}
\usepackage{xcolor}
\usepackage{listings}
\usepackage{fancyhdr}
\usepackage{hyperref}
\usepackage{float}
\usepackage{caption}
\usepackage{subcaption}
\usepackage{tikz}
\usetikzlibrary{shapes.geometric, arrows.meta, positioning}

% --- COLORES CORPORATIVOS Y CÓDIGO ---
\definecolor{primaryBlue}{RGB}{24, 76, 120}
\definecolor{accentOrange}{RGB}{230, 81, 0}
\definecolor{darkGreen}{RGB}{46, 125, 50}
\definecolor{lightGray}{RGB}{245, 247, 250}
\definecolor{codeBackground}{RGB}{248, 249, 250}
\definecolor{codeKeyword}{RGB}{0, 51, 179}
\definecolor{codeString}{RGB}{6, 125, 23}
\definecolor{codeComment}{RGB}{128, 128, 128}

% --- CONFIGURACIÓN DE LISTINGS (KOTLIN) ---
\lstdefinelanguage{Kotlin}{
  keywords={package, import, val, var, fun, class, interface, data, object, private, override, by, lazy, when, if, else, return, init},
  keywordstyle=\color{codeKeyword}\bfseries,
  ndkeywords={Boolean, Int, String, Double, Float, List, StateFlow, MutableStateFlow, ViewModel},
  ndkeywordstyle=\color{primaryBlue}\bfseries,
  sensitive=true,
  comment=[l]{//},
  morecomment=[s]{/*}{*/},
  commentstyle=\color{codeComment}\itshape,
  stringstyle=\color{codeString},
  morestring=[b]",
  morestring=[s]{"""*}{*"""}
}

\lstset{
  language=Kotlin,
  backgroundcolor=\color{codeBackground},
  basicstyle=\ttfamily\footnotesize,
  breaklines=true,
  frame=single,
  rulecolor=\color{black!20},
  numbers=left,
  numberstyle=\tiny\color{black!50},
  showspaces=false,
  showstringspaces=false,
  captionpos=b,
  tabsize=2
}

% --- CONFIGURACIÓN DE ENCABEZADOS Y PIES DE PÁGINA ---
\pagestyle{fancy}
\fancyhf{}
\fancyhead[L]{\footnotesize \textsc{UP Chiapas} | Desarrollo de Aplicaciones Móviles}
\fancyhead[R]{\footnotesize \textbf{RawTune Mobile} - Reporte Técnico}
\fancyfoot[C]{\thepage}
\renewcommand{\headrulewidth}{0.5pt}

\hypersetup{
    colorlinks=true,
    linkcolor=primaryBlue,
    urlcolor=accentOrange,
    citecolor=darkGreen
}

% =========================================================================
%                               DOCUMENTO
% =========================================================================
\begin{document}

% --- PORTADA FORMAL ACADÉMICA ---
\begin{titlepage}
    \centering
    \vspace*{1cm}
    
    {\scshape\Large Universidad Politécnica de Chiapas \par}
    \vspace{0.3cm}
    {\scshape\large Ingeniería en Tecnologías de la Información e Innovación Digital \par}
    
    \vspace{1.5cm}
    \rule{\linewidth}{1.5pt} \\[0.4cm]
    {\huge \bfseries RawTune Mobile: Optimizador y Simulador de Hardware PC/Mobile \\[0.3cm]}
    {\large \textbf{Reporte Arquitectónico, MVVM y Persistencia Local (Unidad II)}}
    \rule{\linewidth}{1.5pt}
    
    \vspace{1.8cm}
    
    \begin{minipage}[t]{0.45\textwidth}
        \textbf{Asignatura:} \\
        Desarrollo de Aplicaciones Móviles \\[0.3cm]
        \textbf{Período Cuatrimestral:} \\
        Septiembre -- Diciembre 2026 \\[0.3cm]
        \textbf{Docente Titular:} \\
        Academia de Desarrollo Móvil
    \end{minipage}
    \hfill
    \begin{minipage}[t]{0.45\textwidth}
        \begin{flushright}
            \textbf{Estudiante:} \\
            Salazar (253448-spec) \\[0.3cm]
            \textbf{Matrícula:} \\
            253448 \\[0.3cm]
            \textbf{Correo Institucional:} \\
            \texttt{253448@it2id.upchiapas.edu.mx}
        \end{flushright}
    \end{minipage}
    
    \vfill
    {\large Suchiapa, Chiapas, México \\ \today \par}
\end{titlepage}

\tableofcontents
\newpage

% =========================================================================
% 1. INTRODUCCIÓN Y OBJETIVOS
% =========================================================================
\section{Introducción y Objetivos del Proyecto}

El presente proyecto documenta el diseño, construcción y validación técnica de \textbf{RawTune Mobile}, una aplicación nativa desarrollada bajo el sistema operativo Android utilizando el toolkit moderno declarativo \textbf{Jetpack Compose} y el lenguaje \textbf{Kotlin}.

\subsection{Objetivo General}
Diseñar e implementar una solución móvil que cumpla con los lineamientos de la \textbf{Unidad II: Diseño de aplicaciones móviles}, aplicando estrictamente la arquitectura por capas \textbf{Model-View-ViewModel (MVVM)}, el patrón de Inyección de Dependencias manual y el almacenamiento local mediante mecanismos de clave-valor (\texttt{SharedPreferences}), respetando la restricción técnica absoluta de \textbf{no emplear bases de datos relacionales (SQLite / Room) ni servicios API remotos}.

\subsection{Dominio de la Solución (RawTune)}
RawTune es una herramienta diagnóstica y optimizadora de especificaciones de hardware. Permite al usuario modelar ensambles computacionales seleccionando procesadores (CPU), capacidades de memoria RAM (de 4 a 64 GB) y tecnologías de bus (DDR4 vs. DDR5). El sistema calcula en tiempo real métricas de ancho de banda teórico en gigabytes por segundo (GB/s), porcentajes de cuello de botella (\textit{bottleneck}) y calificaciones integrales de rendimiento, ofreciendo retroalimentación contextual y recomendaciones técnicas sobre memorias de doble canal y gráficos integrados (iGPU).

% =========================================================================
% 2. FASE 1: ESPECIFICACIÓN DE REQUISITOS (ENTREGABLE 1)
% =========================================================================
\section{Fase 1: Especificación de Requisitos}

\subsection{Requisitos Funcionales (RF)}
\begin{table}[H]
\centering
\small
\begin{tabular}{@{}llp{10cm}@{}}
\toprule
\textbf{ID} & \textbf{Nombre} & \textbf{Descripción} \\ \midrule
\textbf{RF-01} & Selección de CPU & Permite al usuario seleccionar un procesador desde un catálogo inicial o personalizado. \\
\textbf{RF-02} & Configuración de RAM & Permite configurar la capacidad de memoria RAM mediante un deslizador interactivo (\textit{Slider}) de 4 a 64 GB. \\
\textbf{RF-03} & Selección de Tecnología & Permite alternar la tecnología de memoria entre DDR4 (3200 MT/s) y DDR5 (6000 MT/s) mediante \textit{FilterChips}. \\
\textbf{RF-04} & Cálculo y Diagnóstico & Procesa de forma concurrente el puntaje global (0-100), ancho de banda y cuello de botella del ensamble. \\
\textbf{RF-05} & Alta Manual de Procesadores & Permite al usuario dar de alta nuevos procesadores a mano con núcleos, hilos, GPU integrado y alertas de RAM. \\
\textbf{RF-06} & Persistencia de Ensambles & Permite guardar configuraciones en el almacenamiento local (\texttt{SharedPreferences}) y recuperarlas posteriormente. \\
\textbf{RF-07} & Eliminación y Gestión & Permite eliminar perfiles guardados o CPUs personalizados creados por el usuario. \\ \bottomrule
\end{tabular}
\caption{Matriz de Requisitos Funcionales del Sistema.}
\end{table}

\subsection{Requisitos No Funcionales (RNF)}
\begin{itemize}
    \item \textbf{RNF-01 (Arquitectura):} Separación estricta de responsabilidades entre \texttt{Presentation}, \texttt{Domain/DI} y \texttt{Data}.
    \item \textbf{RNF-02 (UDF e Inmutabilidad):} El flujo de datos debe ser estrictamente unidireccional (\textit{Unidirectional Data Flow}) empleando \texttt{StateFlow} inmutable.
    \item \textbf{RNF-03 (Stateless Components):} Los componentes de interfaz deben ser desacoplados mediante \textit{State Hoisting}.
    \item \textbf{RNF-04 (Restricción de Persistencia):} Almacenamiento local mediante \texttt{SharedPreferences} estructurado, sin bibliotecas relacionales (Room/SQLite).
    \item \textbf{RNF-05 (Compatibilidad):} Compatibilidad mínima con Android 8.0 (API nivel 26) y compilación nativa en JVM 21 con Kotlin 2.4.
\end{itemize}

% =========================================================================
% 3. FASE 2: DISEÑO ARQUITECTÓNICO Y FLUJO DE DATOS
% =========================================================================
\section{Fase 2: Arquitectura de Software y Patrón MVVM}

La aplicación implementa una arquitectura modular inspirada en \textbf{Clean Architecture} y las directrices de arquitectura oficial de Google Android.

\subsection{Diagrama de Flujo Arquitectónico}
A continuación se ilustra la interacción entre las capas:

\begin{figure}[H]
\centering
\begin{tikzpicture}[
    box/.style={rectangle, draw=primaryBlue, fill=lightGray, rounded corners=3pt, thick, align=center, inner sep=6pt},
    dataBox/.style={rectangle, draw=darkGreen, fill=green!5, rounded corners=3pt, thick, align=center, inner sep=6pt},
    diBox/.style={rectangle, draw=accentOrange, fill=orange!5, rounded corners=3pt, thick, align=center, inner sep=6pt},
    arrow/.style={-{Latex[scale=1.0]}, thick}
]
    % Capa Presentación
    \node[box, fill=blue!10] (page) {\textbf{RawTunePage (Stateful)}\\ \small Conecta ViewModel y ciclo de vida};
    \node[box, below=1.2cm of page] (content) {\textbf{RawTuneContent (Stateless)}\\ \small Renderiza componentes hijos y emite lambdas};
    
    % ViewModel
    \node[box, right=2.0cm of page, fill=blue!15] (vm) {\textbf{RawTuneViewModel}\\ \small Expone StateFlow y ejecuta lógica};
    
    % DI
    \node[diBox, above=1.2cm of vm] (di) {\textbf{RawTuneModule / Factory}\\ \small Inyección manual (\texttt{by lazy})};
    
    % Capa de Datos
    \node[dataBox, right=2.0cm of vm] (repo) {\textbf{RawTuneRepository}\\ \small Interfaz desacoplada};
    \node[dataBox, below=1.0cm of repo] (ds) {\textbf{RawTuneLocalDataSource}\\ \small Catálogo semilla + Fórmulas};
    \node[dataBox, below=1.0cm of ds, fill=yellow!10] (sp) {\textbf{SharedPreferences}\\ \small Archivo clave-valor \texttt{rawtune\_storage}};

    % Conexiones
    \draw[arrow] (page) -- node[left, font=\footnotesize] {UiState inmutable} (content);
    \draw[arrow] (content.west) to[out=150,in=210] node[left, font=\footnotesize] {Lambdas: \texttt{() -> Unit}} (page.west);
    \draw[arrow] (page) -- node[above, font=\footnotesize] {Observa \texttt{uiState}} (vm);
    \draw[arrow] (vm) -- node[above, font=\footnotesize] {Consume} (repo);
    \draw[arrow] (di) -- node[right, font=\footnotesize] {Inyecta} (vm);
    \draw[arrow] (repo) -- node[right, font=\footnotesize] {Delega en} (ds);
    \draw[arrow] (ds) -- node[right, font=\footnotesize] {Lectura / Escritura} (sp);
\end{tikzpicture}
\caption{Diagrama de Flujo Unidireccional de Datos (UDF) y Capas de la Aplicación.}
\label{fig:arquitectura}
\end{figure}

\subsection{Estructura de Paquetes en el Proyecto}
Para satisfacer los estándares de calidad de código y modularidad exigidos por la plantilla del docente, se estructuró el paquete \texttt{com.upchiapas.kt\_template.rawtune}:

\begin{itemize}
    \item \texttt{rawtune.data.model}: Entidades de datos inmutables (\texttt{CpuOption}, \texttt{HardwareSpecs}, \texttt{AnalysisResult}, \texttt{SavedHardwareProfile}).
    \item \texttt{rawtune.data.datasource}: Manejo directo de colecciones en memoria, datos semilla y persistencia en \texttt{SharedPreferences} (\texttt{RawTuneLocalDataSource}).
    \item \texttt{rawtune.data.repository}: Contrato (\texttt{RawTuneRepository}) e implementación concreta (\texttt{RawTuneRepositoryImpl}).
    \item \texttt{rawtune.di}: Módulo de dependencias perezosas (\texttt{RawTuneModule}) enlazado al \texttt{KtApplication}.
    \item \texttt{rawtune.presentation.viewmodels}: \texttt{RawTuneViewModel} y \texttt{RawTuneViewModelFactory}.
    \item \texttt{rawtune.presentation.screens}: \texttt{RawTunePage} (Stateful/Stateless) y \texttt{RawTuneUiState}.
    \item \texttt{rawtune.presentation.components}: 6+ componentes reutilizables sin estado (\texttt{CpuDropdownSelector}, \texttt{RamConfigCard}, \texttt{DiagnosticMetricsCard}, \texttt{AddCpuDialog}, \texttt{SaveProfileSection}, \texttt{SavedProfilesList}).
\end{itemize}

% =========================================================================
% 4. IMPLEMENTACIÓN TÉCNICA Y PERSISTENCIA LOCAL
% =========================================================================
\section{Implementación Técnica y Persistencia Local}

\subsection{Gestión de Estados Reactivos con StateFlow}
El patrón de diseño garantiza inmutabilidad estricta. El estado completo de la pantalla se centraliza en la data class \texttt{RawTuneUiState}, expuesta mediante el patrón \textit{Backing Property}:

\begin{lstlisting}[caption={Patrón Backing Property en RawTuneViewModel.kt}]
class RawTuneViewModel(
    private val repository: RawTuneRepository
) : ViewModel() {

    // Flujo mutable privado para modificaciones internas atómicas
    private val _uiState = MutableStateFlow(RawTuneUiState())
    
    // Flujo público inmutable de solo lectura para la vista
    val uiState: StateFlow<RawTuneUiState> = _uiState.asStateFlow()

    fun updateRam(ramGb: Int) {
        _uiState.update { it.copy(ramGb = ramGb) }
        analyzeHardware()
    }
}
\end{lstlisting}

\subsection{Persistencia Clave-Valor sin SQLite/Room}
En cumplimiento con la restricción técnica de la Unidad II, la persistencia de perfiles y de nuevos procesadores creados por el usuario se gestiona a través de la interfaz nativa \texttt{SharedPreferences}:

\begin{lstlisting}[caption={Persistencia de procesadores en RawTuneLocalDataSource.kt}]
fun saveCustomCpu(cpu: CpuOption) {
    val sp = prefs ?: return
    val rawIds = sp.getString("custom_cpu_ids_list", "") ?: ""
    val existingIds = rawIds.split(",").filter { it.isNotBlank() }.toMutableList()
    
    if (!existingIds.contains(cpu.id)) {
        existingIds.add(0, cpu.id)
    }

    sp.edit()
        .putString("custom_cpu_ids_list", existingIds.joinToString(","))
        .putString("cpu_${cpu.id}_name", cpu.name)
        .putInt("cpu_${cpu.id}_cores", cpu.cores)
        .putInt("cpu_${cpu.id}_threads", cpu.threads)
        .putInt("cpu_${cpu.id}_score", cpu.baseScore)
        .putString("cpu_${cpu.id}_ram_type", cpu.recommendedRamType)
        .putBoolean("cpu_${cpu.id}_has_gpu", cpu.hasIntegratedGpu)
        .putString("cpu_${cpu.id}_gpu_name", cpu.integratedGpuName)
        .putString("cpu_${cpu.id}_desc", cpu.description)
        .putString("cpu_${cpu.id}_ram_alert", cpu.ramRecommendation)
        .apply()
}
\end{lstlisting}

% =========================================================================
% 5. APARTADO DE CAPTURAS DE PANTALLA DE LA APLICACIÓN
% =========================================================================
\section{Fase 3: Resultados Visuales y Galería de la Aplicación}

En esta sección se evidencian las vistas y componentes construidos con Jetpack Compose. 

\begin{figure}[H]
    \centering
    % Para colocar tu imagen en Overleaf:
    % 1. Sube tu imagen en el panel izquierdo (ej. 'pantalla_principal.png')
    % 2. Descomenta la siguiente línea y elimina el placeholder:
    % \includegraphics[width=0.48\textwidth]{pantalla_principal.png}
    \fbox{\parbox[c][8cm][c]{0.45\textwidth}{\centering \textbf{[CAPTURA 1: PANTALLA PRINCIPAL]}\\[0.3cm] \small Inserte aquí captura con el TopAppBar, Selector de CPU, Slider de RAM y Chips DDR4/DDR5.}}
    \hfill
    % \includegraphics[width=0.48\textwidth]{diagnostico_metricas.png}
    \fbox{\parbox[c][8cm][c]{0.45\textwidth}{\centering \textbf{[CAPTURA 2: DIAGNÓSTICO Y MÉTRICAS]}\\[0.3cm] \small Inserte aquí captura mostrando la barra de progreso lineal de puntaje, ancho de banda y alertas de RAM.}}
    \caption{Interfaz Principal: Configuración de Hardware y Diagnóstico en Tiempo Real.}
    \label{fig:interfaz_principal}
\end{figure}

\vspace{0.5cm}

\begin{figure}[H]
    \centering
    % \includegraphics[width=0.48\textwidth]{dialogo_agregar_cpu.png}
    \fbox{\parbox[c][8cm][c]{0.45\textwidth}{\centering \textbf{[CAPTURA 3: ALTA MANUAL DE CPU]}\\[0.3cm] \small Inserte aquí captura del diálogo modal \texttt{AddCpuDialog} mostrando los campos de núcleos, switch de GPU integrado y alertas.}}
    \hfill
    % \includegraphics[width=0.48\textwidth]{historial_persistente.png}
    \fbox{\parbox[c][8cm][c]{0.45\textwidth}{\centering \textbf{[CAPTURA 4: HISTORIAL PERSISTENTE]}\\[0.3cm] \small Inserte aquí captura de la sección inferior con la lista de ensambles guardados en \texttt{SharedPreferences}.}}
    \caption{Funciones Avanzadas: Registro Manual de Procesadores y Persistencia Local.}
    \label{fig:interfaz_avanzada}
\end{figure}

% =========================================================================
% 6. HISTORIAL DE COMMITS
% =========================================================================
\section{Historial Incremental de Commits en Git}

En apego riguroso a la condición obligatoria de contar con un historial evolutivo verificable, el repositorio en GitHub (\url{https://github.com/253448-spec/RawTune}) cuenta con 6 commits secuenciales:

\begin{table}[H]
\centering
\small
\begin{tabular}{@{}cclp{7.5cm}@{}}
\toprule
\textbf{\#} & \textbf{Hash} & \textbf{Etiqueta del Commit} & \textbf{Aporte Técnico Principal} \\ \midrule
1 & \texttt{8a728dd} & Commit 1: Vista RawTuneScreen inicial & Maquetado preliminar en Compose y ViewModel base. \\
2 & \texttt{8f15e75} & Commit 2: Capa de datos y modelos & Creación de modelos de datos, \texttt{RawTuneLocalDataSource} y contrato de repositorio. \\
3 & \texttt{fc90a2f} & Commit 3: Inyección de Dependencias & Creación de \texttt{KtApplication.kt}, \texttt{RawTuneModule} y \texttt{ViewModelFactory}. \\
4 & \texttt{1b789a8} & Commit 4: Persistencia y UiState & Desacoplamiento de \texttt{RawTuneUiState}, persistencia en \texttt{SharedPreferences} y reactividad UDF. \\
5 & \texttt{d18ff1e} & Commit 5: Componentes modulares & 100\% Stateless components, separación Stateful/Stateless y enlace con \texttt{MainActivity}. \\
6 & \texttt{370b5a7} & Commit 6: Alta manual de procesadores & Modal interactivo para agregar CPUs a mano con GPU integrado, alertas y recálculo dinámico. \\ \bottomrule
\end{tabular}
\caption{Historial Oficial de Commits en la Rama \texttt{main}.}
\end{table}

% =========================================================================
% 7. CONCLUSIONES
% =========================================================================
\section{Conclusiones}
El desarrollo del proyecto \textbf{RawTune Mobile} demostró la viabilidad y solidez del patrón de arquitectura MVVM combinado con Jetpack Compose. Se logró una separación absoluta entre las reglas de cálculo de hardware y la interfaz de usuario, garantizando componentes altamente reutilizables y fácilmente testeables gracias al desacoplamiento \textit{Stateless}. 

Asimismo, se validó la persistencia local eficiente mediante \texttt{SharedPreferences}, permitiendo la permanencia de ensambles y la creación dinámica de nuevos componentes sin requerir el sobrecoste de motores SQL ni dependencias externas de red.

% =========================================================================
% 8. REFERENCIAS BIBLIOGRÁFICAS
% =========================================================================
\begin{thebibliography}{99}
\bibitem{compose} Google Developers. (2026). \textit{Jetpack Compose: Guía de Arquitectura de la Interfaz de Usuario y State Hoisting}. Android Open Source Project. \url{https://developer.android.com/jetpack/compose}
\bibitem{clean} Martin, R. C. (2018). \textit{Clean Architecture: A Craftsman's Guide to Software Structure and Design}. Prentice Hall.
\bibitem{stateflow} JetBrains. (2026). \textit{StateFlow and SharedFlow: Reactive Stream Processing in Kotlin Coroutines}. \url{https://kotlinlang.org/api/kotlinx.coroutines/}
\end{thebibliography}

\end{document}
